\scanpage{22}
% L1–2
donc $|D_{*}|$ est aussi la réalisation topologique de l'ens quasi-simplicial $D_*$
% L3–5
L'ens ordonné $D_*$ a la propriété suivante
\textcircled{a} $\forall d\in D_*$, l'ens \inserted{ord} $D_*(d)=\{\delta\in D_*\mid\delta\leq d\}$ est isomorphe à un $\mathscr P^*(I_d)$, pour $I_d$ ens fini
% L6–9
(déterminé à isom can près comme l'ens des él. minimaux de $D_*(d)$, ou encore l'ens des éléments minimaux de $D_*$ -- \deleted{\unsure{qui}} i.e. des él. de $D_0$ -- majorés par $d$)

% L10–14
Mais comment exprimer l'ordre total des $I_d$ \inserted{$(d\in D_*)$} en termes d'une structure suppl. sur l'ens ordonné $D_*$ ? [Souvent elle sera induite par une structure d'ordre (pas total) sur $D_0$, ens des él. minimaux de $D_*$.]
% L15–18
La donnée de cet ordre total, avec la condition que pour $d'<d$, l'ordre total de $I_{d'}$ est induit par celui de $I_d$, \uline{équivaut} à la donnée de l'ens quasi-simplicial $\unsure{\mathbb D|D}_*$. On va examiner des
% L19–21
\margin{$\vert\,\vert\,\vert$}conditions \inserted{sur $D_*$} moyennant les quelles les ordres totaux des $I_d$ $(d\in D_*)$ sont bien induits par un ordre unique sur $D_0$.

% L22–23
Considérons
\diagram{\begin{tikzcd}[ampersand replacement=\&,column sep=large]
\mathbb D_0 \& D_1\arrow[l,shift right=1.5,"\partial_0"']\arrow[l,shift left=1.5,"\partial_1"] \&
D_2\arrow[l,shift right=3,"\partial_0"']\arrow[l,"\partial_1" description]\arrow[l,shift left=3,"\partial_2"]
\end{tikzcd}}
% L24
Considérons Sur $\unsure{D|D_0|\underline D}$ la relation
% archive sheet number: 21
